\documentclass[12pt, a4paper]{article}
\usepackage{../../../myconfig_line} % Gọi file cấu hình từ ngoài gốc vào
% ==========================================
% BẮT ĐẦU NỘI DUNG TÀI LIỆU
% ==========================================
\begin{document}

% Truyền 3 tham số: {Nhãn cam} {Chữ to chính giữa} {Chữ ở góc phải Header}
\titlebox{Chủ đề: Tích phân}{Tốc dộ thay đổi đại lượng và lượng tích lũy}{Thay đổi đại lượng}


% --- CÂU 1 ---
\questionTrueFalse{20}{1}{Bạn Khánh đang tham gia một thử nghiệm nhằm đo thời gian ghi nhớ các khái niệm khác nhau trong môn Nguyên lý kế toán. Biết rằng hàm số \(M(t)\) biểu thị số lượng khái niệm mà bạn Khánh nhớ được trong \(t\) phút, khi đó tốc độ ghi nhớ khái niệm của bạn Khánh được biểu thị bởi hàm số \(M'(t) = 0,2t - 0,005t^2\) (khái niệm/phút). Giả sử rằng bạn Khánh rất yêu môn học này, nên trước khi tham gia thử nghiệm, bạn không thuộc được khái niệm nào.}{
    \textbf{a)} & Họ nguyên hàm của hàm số \(M'(t)\) là \(M(t) = 0,1t^2 - \dfrac{1}{600}t^3 + C\).
    & \textcolor{amberorange}{$\bigcirc$} & \textcolor{amberorange}{$\bigcirc$} \\ \hline
    
    \textbf{b)} & Sau 10 phút, bạn Khánh học thuộc được 9 khái niệm (kết quả làm tròn đến hàng đơn vị).
    & \textcolor{amberorange}{$\bigcirc$} & \textcolor{amberorange}{$\bigcirc$} \\ \hline
    
    \textbf{c)} & Từ phút thứ 15 đến phút thứ 30, bạn Khánh học thuộc được 28 khái niệm (kết quả làm tròn đến hàng đơn vị).
    & \textcolor{amberorange}{$\bigcirc$} & \textcolor{amberorange}{$\bigcirc$} \\ \hline
    
    \textbf{d)} & Nếu thử nghiệm này kéo dài hơn 40 phút, bạn Khánh sẽ bắt đầu quên dần các khái niệm đã học (tức số khái niệm học thuộc được giảm dần).
    & \textcolor{amberorange}{$\bigcirc$} & \textcolor{amberorange}{$\bigcirc$} \\
}

% --- CÂU 2 ---
\questionShortAnswer{30}{2}{Một công ty hóa chất có một bồn chứa dạng hình nón cụt làm bằng thép, bồn có chiều cao là 4 mét, bán kính đáy dưới là 1 mét và bán kính đáy trên là 3 mét. Giả sử bồn đang trống, người ta bắt đầu bơm một loại dung dịch vào bồn với tốc độ không đổi là \(0,5\) m\(^3\)/phút. Hỏi sau bao lâu (phút) kể từ khi bắt đầu bơm, mực nước trong bồn đạt độ cao 2 mét? (kết quả đến chữ số hàng phần mười).
}

% --- CÂU 3 ---
\questionShortAnswer{30}{3}{Những ngày giáp Tết Nguyên Đán cũng là dịp bước vào vụ Đông Xuân, cây lúa sau khi được cấy trải qua quá trình tăng trưởng để nhánh và phát triển chiều cao trước khi làm đòng, trổ bông. Qua nghiên cứu một giống lúa mới, các nhà khoa học nhận thấy một cây lúa tính từ lúc được cấy bằng một cây mạ với chiều cao 15 cm có tốc độ tăng trưởng chiều cao cho bởi hàm số \(v(t) = -0,3t^3 + 2,2t^2\), trong đó \(t\) tính theo tuần, \(v(t)\) tính bằng cm/tuần. Gọi \(h(t)\) là chiều cao của cây lúa ở tuần thứ \(t\) (\(t \geq 0, t \in \mathbb{R}\)). Biết rằng khi đạt chiều cao lớn nhất cây lúa sẽ dừng phát triển chiều cao để tập trung dinh dưỡng nuôi bông. Tính chiều cao lớn nhất của cây lúa (đơn vị cm, kết quả làm tròn đến hàng đơn vị).
}

% --- CÂU 4 ---
\questionShortAnswer{30}{4}{Tốc độ giải ngân 2 tỷ tiền trợ cấp dành cho một vùng bị thiệt hại lũ lụt tỉ lệ thuận với bình phương của \((100-t)\), trong đó \(t\) là thời gian tính bằng ngày, \(0 \leq t \leq 100\), và \(M(t)\) là số tiền còn lại chưa giải ngân sau \(t\) ngày. Hỏi số tiền còn lại chưa giải ngân sau 40 ngày là bao nhiêu triệu đồng, biết rằng toàn bộ số tiền sẽ được giải ngân trong 100 ngày.
}

\end{document}